\subparagraph{Proofs - Separoid Axioms for Conditional Independence}\hfill%

In the following let $(\Wcal \times \Tcal, K(W|T))$ be a transition probability space with Markov kernel:
\[K(W|T): \, \Tcal \dashrightarrow  \Wcal,\]
and \trv{}s $X$, $Y$, $Z$, $U$ with common domain $\Wcal\times \Tcal$ and measurable spaces $\Xcal$, 
$\Ycal$, $\Zcal$, $\Ucal$, resp., as codomains. We indicate when we need to assume standard measurable spaces.\\
We will use
$T:\, \Wcal \times \Tcal \to \Tcal$ to denote the canonical projection and $\ast$ to denote the constant \trv{}.\\
Recall that we write $U \precsim X$ if there exists a measurable function  
$\varphi:\, \Xcal \to \Ucal$  such that $U = \varphi \circ X$.\\
Recall that for proving:
\[X \Indep_{K(W|T)} Y \given Z, \]
we need to find/construct a Markov kernel $Q(X|Z)$ such that:
\[ K(X,Y,Z|T) = Q(X|Z) \otimes K(Y,Z|T),\]
which is equivalent to:\\ 
For all $t \in \Tcal$ and all measurable $A \ins \Xcal$, $B \ins \Ycal$, $C \ins \Zcal$ we have the equation:
\begin{align*}
    \E_t\lB Q(X \in A|Z_t) \cdot \I_B(Y_t) \cdot \I_C(Z_t)\rB &= \E_t\lB\I_A(X_t) \cdot \I_B(Y_t) \cdot \I_C(Z_t) \rB 
\end{align*}
where the expectation value $\E_t$ is w.r.t.\ $K(W|T=t)$.\\
We abbreviate $\Indep := \Indep_{K(W|T)}$ in the following.


\begin{Lem}[Extended Left Redundancy]
    \[U \precsim X \implies U \Indep  Y \given X.\]
\begin{proof}
    If $U=\varphi(X)$ put
    $Q(U \in D|X=x):=\delta_\varphi(U \in D|X=x):=\I_D(\varphi(x))$. Then we get:
\begin{align*} 
    \E_t\lB \delta_\varphi(U \in D|X_t) \cdot \I_B(Y_t) \cdot \I_C(X_t) \rB &= \E_t\lB \I_D(U_t) \cdot \I_B(Y_t) \cdot \I_C(X_t) \rB .     
\end{align*}
In suggestive symbols:
\[ K(U|\cancel{T,Y},X) = \delta_\varphi(U|X). \]
\end{proof}
\end{Lem}

\begin{Lem}[$T$-Restricted Right Redundancy] Let $\Xcal$ and $\Zcal$ be standard measurable spaces. Then:
    \[X \Indep  \ast \given Z,T \qquad \text{holds.}\]
\begin{proof}
    Because $\Xcal$ and $\Zcal$ are standard measurable spaces
    we have a factorization:
    \[ K(X,\ast, Z,T|T) = K(X|Z,T) \otimes K(\ast, Z,T|T).\]
    with the conditional Markov kernel $K(X|Z,T)$ of $K(X,Z|T)$ (via theorem \ref{thm-regular-conditional-Markov-kernel}).
    This already shows the claim. In suggestive symbols:
    \[K(X|\cancel{\ast,T},Z,T) = K(X|Z,T).\]
\end{proof}
\end{Lem}

\begin{Lem}[$T$-Inverted Right Decomposition]
    \[X \Indep  Y \given Z \implies X \Indep  T, Y \given Z.\]
\begin{comment}
\begin{proof}
    We use the same $Q(X|Z)$. Then we get by assumption indicated by (!):
    \begin{align*}                  
      & \quad\; \E_t\lB Q( X \in A| Z_t) \cdot\I_B(Y_t) \cdot \I_{C}(Z_t) \cdot \I_D(T_t) \cdot \I_E(Z_t)\rB \\
      & = \E_t\lB Q( X \in A| Z_t) \cdot\I_B(Y_t) \cdot \I_{C}(Z_t) \cdot \I_D(t) \cdot \I_E(Z_t)\rB  \\
      & = \E_t\lB Q( X \in A| Z_t) \cdot\I_B(Y_t) \cdot \I_{C \cap E}(Z_t) \rB \cdot \I_D(t)  \\
      & \stackrel{(!)}{=}  \E_t\lB \I_A(X_t) \cdot\I_B(Y_t) \cdot \I_{C \cap E}(Z_t) \rB \cdot \I_D(t) \\
      & = \E_t\lB \I_A(X_t) \cdot\I_B(Y_t) \cdot \I_{C \cap E}(Z_t) \cdot \I_D(t)\rB \\
      & =  \E_t\lB \I_A(X_t) \cdot\I_B(Y_t) \cdot \I_C(Z_t) \cdot \I_D(T_t) \cdot \I_E(Z_t)\rB.
    \end{align*}
In suggestive symbols:
\[ K(X|\cancel{T,T,Y,Z},Z) = K(X|\cancel{T,Y},Z).  \]
\end{proof}
\end{comment}
\begin{proof}
    We use the same $Q(X|Z)$. Then we get by assumption indicated by (!):
    \begin{align*}                  
      & \quad\; \E_t\lB Q( X \in A| Z_t) \cdot\I_B(Y_t)  \cdot \I_D(T_t) \cdot \I_C(Z_t)\rB \\
      & = \E_t\lB Q( X \in A| Z_t) \cdot\I_B(Y_t)  \cdot \I_D(t) \cdot \I_C(Z_t)\rB  \\
      & = \E_t\lB Q( X \in A| Z_t) \cdot\I_B(Y_t) \cdot \I_C(Z_t) \rB \cdot \I_D(t)  \\
      & \stackrel{(!)}{=}  \E_t\lB \I_A(X_t) \cdot\I_B(Y_t) \cdot \I_C(Z_t) \rB \cdot \I_D(t) \\
      & = \E_t\lB \I_A(X_t) \cdot\I_B(Y_t) \cdot \I_C(Z_t) \cdot \I_D(t)\rB \\
      & =  \E_t\lB \I_A(X_t) \cdot\I_B(Y_t) \cdot \I_D(T_t) \cdot \I_C(Z_t)\rB.
    \end{align*}
In suggestive symbols:
\[ K(X|\cancel{T,T,Y},Z) = K(X|\cancel{T,Y},Z).  \]
\end{proof}
\end{Lem}

\begin{comment}
\begin{Lem}[Left Decomposition]
    \[(X \Indep  Y \given Z) \land (U \precsim X) \implies U \Indep  Y \given Z.\]
\begin{proof}
    Let $Q(X|Z)$ be given from the left and $U =\varphi(X)$. Then we use the push-forward Markov kernel $Q(\varphi(X)|Z)$.
    We then get by assumption (!):
    \begin{align*}                  
      & \quad\; \E_t\lB Q( \varphi(X) \in D| Z_t) \cdot\I_B(Y_t) \cdot \I_C(Z_t) \rB \\
      & = \E_t\lB Q( X \in \varphi^{-1}(D)| Z_t) \cdot\I_B(Y_t) \cdot \I_C(Z_t) \rB \\
      & \stackrel{(!)}{=} \E_t\lB \I_{\varphi^{-1}(D)}(X_t) \cdot\I_B(Y_t) \cdot \I_C(Z_t) \rB \\
      & = \E_t\lB \I_D(\varphi(X_t)) \cdot\I_B(Y_t) \cdot \I_{C}(Z_t) \rB\\
      & = \E_t\lB \I_D(U_t) \cdot\I_B(Y_t) \cdot \I_{C}(Z_t) \rB.
    \end{align*}
 In suggestive symbols:
 \[ K(U|\cancel{T,Y},Z) = K(\varphi(X)|\cancel{T,Y},Z). \]
\end{proof}
\end{Lem}
\end{comment}

\begin{Lem}[Left Decomposition]
    \[X,U \Indep  Y \given Z \implies U \Indep  Y \given Z.\]
\begin{proof}
    Let $Q(X,U|Z)$ be given from the left. Then we use the marginal Markov kernel $Q(U|Z)$.
    We then get by assumption (!):
    \begin{align*}                  
      & \quad\; \E_t\lB Q(U \in D| Z_t) \cdot\I_B(Y_t) \cdot \I_C(Z_t) \rB \\
      & = \E_t\lB Q( X \in \Xcal, U \in D| Z_t) \cdot\I_B(Y_t) \cdot \I_C(Z_t) \rB \\
      & \stackrel{(!)}{=} \E_t\lB \I_\Xcal(X_t) \cdot \I_D(U_t) \cdot\I_B(Y_t) \cdot \I_C(Z_t) \rB \\
      & = \E_t\lB \I_D(U_t) \cdot\I_B(Y_t) \cdot \I_{C}(Z_t) \rB.
    \end{align*}
 In suggestive symbols:
 \[ K(U|\cancel{T,Y},Z) = K(X \in \Xcal,U)|\cancel{T,Y},Z). \]
\end{proof}
\end{Lem}

\begin{comment}
\begin{Lem}[Right Decomposition]
    \[(X \Indep  Y \given Z) \land (U\precsim Y) \implies X \Indep  U \given Z.\]
\begin{proof}
    Let $Q(X|Z)$ be given from the left and $U=\varphi(Y)$. We then have by assumption (!):
    \begin{align*}                  
      & \quad\; \E_t\lB Q( X \in A| Z_t) \cdot\I_D(U_t) \cdot \I_C(Z_t) \rB \\
      & = \E_t\lB Q( X \in A| Z_t) \cdot\I_D(\varphi(Y_t)) \cdot \I_C(Z_t) \rB \\
      & = \E_t\lB Q( X \in A| Z_t) \cdot\I_{\varphi^{-1}(D)}(Y_t) \cdot \I_C(Z_t) \rB \\
      & \stackrel{(!)}{=} \E_t\lB \I_A(X_t) \cdot\I_{\varphi^{-1}(D)}(Y_t)  \cdot \I_C(Z_t) \rB \\
      & = \E_t\lB \I_A(X_t) \cdot\I_D(\varphi(Y_t)) \cdot \I_{C}(Z_t) \rB\\
      & = \E_t\lB \I_A(X_t) \cdot\I_D(U_t) \cdot \I_{C}(Z_t) \rB.
    \end{align*}
In suggestive symbols:
\[ K(X|\cancel{T,\varphi(Y)},Z) = K(X|\cancel{T,Y},Z). \]
\end{proof}
\end{Lem}
\end{comment}

\begin{Lem}[Right Decomposition]
    \[X \Indep  Y,U \given Z \implies X \Indep  U \given Z.\]
\begin{proof}
    Let $Q(X|Z)$ be given from the left. We then have by assumption (!):
    \begin{align*}                  
      & \quad\; \E_t\lB Q( X \in A| Z_t) \cdot\I_D(U_t) \cdot \I_C(Z_t) \rB \\
      & = \E_t\lB Q( X \in A| Z_t) \cdot \I_\Ycal(Y_t) \cdot\I_D(U_t) \cdot \I_C(Z_t) \rB \\
      & \stackrel{(!)}{=} \E_t\lB \I_A(X_t) \cdot \I_\Ycal(Y_t) \cdot\I_D(U_t)  \cdot \I_C(Z_t) \rB \\
      & = \E_t\lB \I_A(X_t) \cdot\I_D(U_t) \cdot \I_{C}(Z_t) \rB.
    \end{align*}
In suggestive symbols:
\[ K(X|\cancel{T,U},Z) = K(X|\cancel{T,Y,U},Z). \]
\end{proof}
\end{Lem}

\begin{Lem}[Left Weak Union] Let $\Xcal$ and $\Ucal$ be standard measurable spaces. Then:
    \[X , U \Indep  Y  \given Z \implies X \Indep  Y \given U,Z.\]
\begin{proof}
    By assumption we have:
    \[K(X,U,Y,Z|T) = Q(X,U|Z) \otimes K(Y,Z|T),\]
    for some Markov kernel $Q(X,U|Z)$. 
    If we marginalize out $X$ we get:
    \[K(U,Y,Z|T) = Q(U|Z) \otimes K(Y,Z|T).\]
    Because $\Xcal$ and $\Ucal$ are standard measurable spaces
    we have a factorization:
    \[ Q(X,U|Z) = Q(X|U,Z) \otimes Q(U|Z).\]
    with the conditional Markov kernel $Q(X|U,Z)$ (via theorem \ref{thm-regular-conditional-Markov-kernel}).\\
    Putting these equations together we get:
    \begin{align*}
        K(X,U,Y,Z|T) &= Q(X,U|Z) \otimes K(Y,Z|T)\\
                     &= Q(X|U,Z) \otimes Q(U|Z) \otimes K(Y,Z|T) \\
                     &= Q(X|U,Z) \otimes  K(U,Y,Z|T).
    \end{align*}
  In suggestive symbols, this means that:
  $K(X|\cancel{T,Y},U,Z)$ is the conditional of $K(X,U|\cancel{T,Y},Z)$.
\end{proof}
\end{Lem}

\begin{Lem}[Right Weak Union] 
    \[X \Indep  Y , U \given Z \implies X \Indep  Y \given U , Z.\]
\begin{proof}
    We have the factorization:
    \[K(X,Y,U,Z|T) = Q(X|Z) \otimes K(Y,U,Z|T),\]
    with some Markov kernel $Q(X|Z)$. If we view $Q(X|Z)$ as a function in $(u,z)$ via:
    \[ (u,z) \mapsto Q(X|Z=z), \]
    by just ignoring the argument $u$ then the claim follows from the same factorization above.\\
    In suggestive symbols:
    \[K(X|\cancel{T,Y},U,Z) = K(X|\cancel{T,Y,U},Z).\]
\end{proof}
\end{Lem}

\begin{Lem}[Left Contraction] 
    \[ (X \Indep  Y \given U , Z) \land (U \Indep  Y \given Z) \implies X , U \Indep  Y \given Z.\]
\begin{proof}
    By assumption we have the two factorizations:
    \begin{align*}
        K(X,Y,U,Z|T) & = Q(X|U,Z) \otimes K(Y,U,Z|T),\\
        K(Y,U,Z|T) & = P(U|Z) \otimes K(Y,Z|T),
    \end{align*}
    with some Markov kernels $Q(X|U,Z)$, $P(U|Z)$. Putting these equations together using $Q(X|U,Z) \otimes P(U|Z)$ we get:
    \[ K(X,Y,U,Z|T) = \lp Q(X|U,Z) \otimes P(U|Z)\rp \otimes K(Y,Z|T).\]
    In suggestive symbols:
    \[K(X,U|\cancel{T,Y},Z) = K(X|\cancel{T,Y},U,Z) \otimes K(U|\cancel{T,Y},Z).\]
\end{proof}
\end{Lem}

\begin{Lem}[Right Contraction]
    \[ (X \Indep  Y \given U , Z) \land (X \Indep  U \given Z)  \implies X \Indep  Y , U \given Z.\]
\begin{proof}
 By assumption we have the two factorizations:
    \begin{align*}
        K(X,Y,U,Z|T) & = Q(X|U,Z) \otimes K(Y,U,Z|T),\\
        K(X,U,Z|T) & = P(X|Z) \otimes K(U,Z|T),
    \end{align*}
    with some Markov kernels $Q(X|U,Z)$, $P(X|Z)$.\\
    Marginalizing out $Y$ we get the equalities:
    \begin{align*}
        K(X,U,Z|T) & = Q(X|U,Z) \otimes K(U,Z|T),\\
        K(X,U,Z|T) & = P(X|Z) \otimes K(U,Z|T).
    \end{align*}
    By the essential uniqueness (see lemma \ref{ess-unique}) of such factorization we get that for every $A \in \Bcal_\Xcal$:
\[ Q(X \in A|U,Z) = P(X \in A|Z) \qquad K(U,Z|T)\text{-a.s.}\]
The same equation then holds also $K(Y,U,Z|T)$-a.s. (by ignoring argument $y$).
Plugging that back into the first equation gives:
\[ K(X,Y,U,Z|T) = P(X|Z) \otimes K(Y,U,Z|T).\]
In suggestive symbols:
\[ K(X|\cancel{T,Y,U},Z) = K(X|\cancel{T,Y},U,Z) =  K(X|\cancel{T,U},Z) \qquad \text{a.s.}\]
\end{proof}
\end{Lem}

\begin{Lem}[Right Cross Contraction]
    \[ (X \Indep  Y \given U , Z) \land (U \Indep  X \given Z) \implies X  \Indep  Y , U \given Z.\]
\begin{proof}
     By assumption we have the two factorizations:
    \begin{align}
        K(X,Y,U,Z|T) & = Q(X|U,Z) \otimes K(Y,U,Z|T), \label{eq:rcc-a}\\
        K(X,U,Z|T) & = P(U|Z) \otimes K(X,Z|T), \label{eq:rcc-b}
    \end{align}
     with some Markov kernels $Q(X|U,Z)$, $P(U|Z)$.\\
     We then define the Markov kernel:
    \begin{align}
        R(X,U|Z) & := Q(X|U,Z) \otimes P(U|Z). \label{eq:rcc-c}
    \end{align}
    We will now show that its marginal: 
    \begin{align}
        R(X|Z) & = Q(X|U,Z) \circ P(U|Z). \label{eq:rcc-c2}
    \end{align} 
    will satisfy the claim.\\
    If we marginalize out $Y$ from equation \ref{eq:rcc-a} we get:
    \begin{align}
        K(X,U,Z|T) & = Q(X|U,Z) \otimes K(U,Z|T). \label{eq:rcc-d}
    \end{align}
    Equating equations \ref{eq:rcc-b} and \ref{eq:rcc-d} gives:
    \begin{align}
        P(U|Z) \otimes K(X,Z|T) &=      K(X,U,Z|T)  
                                 = Q(X|U,Z) \otimes K(U,Z|T). \label{eq:rcc-q}
    \end{align}
    Marginalizing out $X$ in equation \ref{eq:rcc-q} on both sides gives:
    \begin{align}
        K(U,Z|T) &=  P(U|Z) \otimes K(Z|T).  \label{eq:rcc-r}
    \end{align}
    If we now plug equation \ref{eq:rcc-r} into \ref{eq:rcc-d} then we get:
    \begin{align}
        K(X,U,Z|T) & = Q(X|U,Z) \otimes P(U|Z) \otimes K(Z|T) \\
            &\stackrel{\ref{eq:rcc-c}}{=} R(X,U|Z) \otimes K(Z|T). \label{eq:rcc-e}
    \end{align}
    If we marginalize out $U$ in equation \ref{eq:rcc-e} and use definition \ref{eq:rcc-c2}  we arrive at:
    \begin{align}
        K(X,Z|T) & = R(X|Z) \otimes K(Z|T). \label{eq:rcc-f}
    \end{align}
    We now get:
    \begin{align}
        Q(X|U,Z) \otimes K(U,Z|T) 
        & \stackrel{\ref{eq:rcc-d}}{=}  K(X,U,Z|T) \\
        & \stackrel{\ref{eq:rcc-b}}{=} P(U|Z) \otimes K(X,Z|T) \\
        &\stackrel{\ref{eq:rcc-f}}{=} P(U|Z) \otimes R(X|Z) \otimes K(Z|T) \\
        &=  R(X|Z)  \otimes P(U|Z) \otimes K(Z|T) \\
        &\stackrel{\ref{eq:rcc-r}}{=} R(X|Z)  \otimes  K(U,Z|T).
    \end{align}
    By the essential uniqueness (see lemma \ref{ess-unique}) of such a factorization we get that for every $A \in \Bcal_\Xcal$:
    \begin{align}
        Q(X \in A|U,Z) &= R(X \in A|Z) \qquad K(U,Z|T)\text{-a.s.}  \label{eq:ess-un}
    \end{align}
The same equation then holds also $K(Y,U,Z|T)$-a.s. (by ignoring the non-occurring argument $y$).
Plugging \ref{eq:ess-un} back into the equation \ref{eq:rcc-a} we get:
    \begin{align}
        K(X,Y,U,Z|T) & = Q(X|U,Z) \otimes K(Y,U,Z|T),\\
                     & = R(X|Z)  \otimes K(Y,U,Z|T).
    \end{align}
This shows the claim.\\
In suggestive symbols:
\[ K(X|\cancel{T,Y,U},Z) = K(X|\cancel{T,Y},U,Z) \circ K(U|\cancel{T,X},Z).\]
\end{proof}
\end{Lem}

\begin{Lem}[Flipped Left Cross Contraction]
    \[ (X \Indep  Y \given U , Z) \land (Y \Indep  U \given Z) \implies Y \Indep  X , U \given Z.\]
\begin{proof}
 By assumption we have the two factorizations:
    \begin{align*}
        K(X,Y,U,Z|T) & = Q(X|U,Z) \otimes K(Y,U,Z|T),\\
        K(Y,U,Z|T) & = P(Y|Z) \otimes K(U,Z|T),
    \end{align*}
     with some Markov kernels $Q(X|U,Z)$, $P(Y|Z)$.\\
     Marginalizing out $Y$ in the first equation we get the equality:
    \begin{align*}
        K(X,U,Z|T) & = Q(X|U,Z) \otimes K(U,Z|T).
    \end{align*}
Plugging all three equations into each other we get:
    \begin{align*}
        K(X,Y,U,Z|T) & = Q(X|U,Z) \otimes K(Y,U,Z|T)\\
                     & = Q(X|U,Z) \otimes  P(Y|Z) \otimes K(U,Z|T)\\
                     & = P(Y|Z) \otimes  Q(X|U,Z)   \otimes K(U,Z|T) \\
                     & =  P(Y|Z)  \otimes   K(X,U,Z|T).
    \end{align*}
In suggestive symbols:
\[ K(Y|\cancel{T,X,U},Z) = K(Y|\cancel{T,U},Z). \]
\end{proof}
\end{Lem}


\begin{comment}
\begin{Lem}[Symmetry]
    Let $\Ycal$ and $\Zcal$ be standard measurable spaces and $\Tcal=\Asterisk=\{\ast\}$ the one-point space. 
    Then:
    \[X \Indep  Y \given Z \implies Y \Indep X \given Z.\]
\begin{proof}
    By assumption we have:
    \[ K(X,Y,Z|\ast) = Q(X|Z) \otimes K(Y,Z|\ast).  \]
    If we marginalize out $Y$ we get:
    \[ K(X,Z|\ast) = Q(X|Z) \otimes K(Z|\ast).\]
Because $\Ycal$ and $\Zcal$ are standard measurable spaces
    we have a factorization:
    \[ K(Y,Z|\ast) = K(Y|Z,\ast) \otimes K(Z|\ast),\]
    with the conditional Markov kernel $K(Y|Z,\ast)$ (via theorem \ref{thm-regular-conditional-Markov-kernel}).\\
 Using the above equations we get:
    \begin{align*}
        K(X,Y,Z|\ast) &= Q(X|Z) \otimes K(Y,Z|\ast)\\
                      &= Q(X|Z) \otimes K(Y|Z,\ast) \otimes K(Z|\ast)\\
                      &= K(Y|Z,\ast) \otimes Q(X|Z) \otimes K(Z|\ast)\\
                      &= K(Y|Z,\ast) \otimes K(X,Z|\ast).
    \end{align*}
    If we put $P(Y|Z) := K(Y|Z,\ast)$ we get the claim.\\
In suggestive symbols:
%$K(Y|\cancel{X},Z)$ is the conditional of $K(Y,Z|\ast)$.
\[ K(Y|\cancel{\ast,X},Z) = K(Y|Z,\ast).\]
\end{proof}
\end{Lem}
\end{comment}

\begin{Lem}[$T$-Restricted Symmetry]
    Let $\Ycal$ and $\Zcal$ be standard measurable spaces. 
    Then:
    \[X \Indep  Y \given Z,T \implies Y \Indep X \given Z,T.\]
\begin{proof}
    This follows from Flipped Left Cross Contraction with $U=\ast$ and $(Z,T)$ for $Z$:
\[ \lp X \Indep Y \given Z,T \rp \quad \land\quad \lp Y \Indep \ast \given Z,T \rp \implies Y \Indep X \given Z,T, \]
together with $T$-Restricted Right Redundancy:
\[ Y \Indep \ast \given Z,T.  \]
In suggestive symbols:
\[ K(Y|\cancel{\ast,X},Z) = K(Y|Z,\ast).\]
\end{proof}
\end{Lem}
